\documentclass[aspectratio=169, 10pt]{beamer}

% --- PAQUETES BÁSICOS ---
\usepackage[utf8]{inputenc}
\usepackage[T1]{fontenc}
\usepackage[spanish, es-tabla]{babel}
\usepackage{amsmath}
\usepackage{amssymb}
\usepackage{booktabs}
\usepackage{tikz}
\usetikzlibrary{shapes.geometric, arrows, positioning}

% --- TEMA FIME UANL ---
\usetheme[showlogo, nosectionpage]{FIME}

% Estilos TikZ adaptados
\tikzset{
    block/.style={rectangle, draw=FIMEgreen, fill=FIMElightgray, text width=2.4cm, text centered, rounded corners, minimum height=1.2cm, font=\scriptsize\bfseries, line width=1pt},
    arrow/.style={thick,->,>=stealth,draw=UANLblue,line width=1.2pt},
    actor/.style={rectangle, draw=UANLblue, fill=FIMElightgray, text width=2.8cm, text centered, rounded corners, minimum height=1cm, font=\small\bfseries, line width=1pt}
}

% --- METADATOS ---
\title[Herramientas Multicriterio - Sesión 11]{Herramientas Multicriterio para Toma de Decisiones}
\subtitle{Sesión 11: Modelos Híbridos: Integración de MCDM con Optimización (NSGA-II)}
\author[Leonardo G. Hernández]{Dr. Leonardo Gabriel Hernández Landa}
\institute[FIME - UANL]{
    Doctorado en Logística y Cadena de Suministro\\
    Facultad de Ingeniería Mecánica y Eléctrica (FIME)\\
    Universidad Autónoma de Nuevo León (UANL)
}
\date{Tetramestre de Otoño}

\begin{document}

% Slide 1: Portada
\begin{frame}[plain]
    \titlepage
\end{frame}

% Slide 2: Objetivos de la Sesión
\begin{frame}{Objetivos de la Sesión}
    \begin{block}{Al finalizar esta sesión, el estudiante de doctorado será capaz de:}
        \begin{itemize}
            \item \textbf{Analizar críticamente} los fundamentos teóricos del tema: Modelos Híbridos: Integración de MCDM con Optimización (NSGA-II).
            \item \textbf{Implementar y resolver} los modelos matemáticos correspondientes en problemas de logística.
            \item \textbf{Evaluar} la robustez de las decisiones resultantes.
        \end{itemize}
    \end{block}
\end{frame}

% Slide 3: Agenda de la Sesión (3 Horas)
\begin{frame}{Agenda de la Sesión (3 Horas)}
    \begin{columns}[T]
        \begin{column}{0.31\textwidth}
            \begin{block}{Bloque 1 (90 min)}
                \textbf{Teoría y Algoritmos}
                \begin{itemize}
                    \scriptsize
                    \item Conceptos clave del tema.
                    \item Formulación matemática del modelo.
                \end{itemize}
            \end{block}
        \end{column}
        
        \begin{column}{0.31\textwidth}
            \begin{exampleblock}{Bloque 2 (45 min)}
                \textbf{Debate de Literatura}
                \begin{itemize}
                    \scriptsize
                    \item Análisis de lecturas asignadas.
                    \item Discusión socrática sobre aplicaciones logísticas.
                \end{itemize}
            \end{exampleblock}
        \end{column}
        
        \begin{column}{0.31\textwidth}
            \begin{alertblock}{Bloque 3 (45 min)}
                \textbf{Taller Aplicado}
                \begin{itemize}
                    \scriptsize
                    \item Práctica guiada en Python/Excel.
                    \item Planteamiento de la actividad autónoma.
                \end{itemize}
            \end{alertblock}
        \end{column}
    \end{columns}
\end{frame}

% Slide 4: Contenido (Plantilla)
\begin{frame}{Fundamentos de Modelos Híbridos: Integración de MCDM con Optimización (NSGA-II)}
    \begin{itemize}
        \item Espacio reservado para el contenido teórico principal de la Sesión 11.
        \item Resumen de fórmulas, algoritmos o diagramas conceptuales.
    \end{itemize}
    \begin{block}{Nota de Investigación}
        El estudiante de doctorado debe relacionar este tema con el avance de su Producto Integrador de Aprendizaje (PIA).
    \end{block}
\end{frame}

\end{document}
